\chapter{Kubernetes 在 AI 集群中的调度}
\label{chap:k8s-ai-scheduling}

前几章我们已经沿着 AI Infrastructure 的物理路径逐层展开：第 17 章讨论单个 GPU 的计算与内存层级，第 18 章讨论单机多卡的 NVLink 与 NVSwitch，第 19 章讨论跨节点的 RDMA、RoCE 与 AI 集群网络。到这里，一个大模型系统的硬件底座已经逐渐清晰：GPU 提供计算，HBM 提供显存，NVLink/NVSwitch 提供节点内高速互联，RDMA/RoCE 提供节点间高速网络。

但是，硬件本身不会自动变成可用的集群能力。一个真实 AI 集群通常同时运行多种任务：

\begin{itemize}
  \item 大规模预训练作业，可能需要几百到几千张 GPU 连续运行数周；
  \item SFT、RLHF、DPO、评测、数据清洗等中小规模离线任务；
  \item 在线 LLM 推理服务，需要低延迟、高可用和弹性伸缩；
  \item embedding、rerank、多模态模型、语音模型等不同形态的服务；
  \item 多个团队、多个项目、多个优先级、多个预算池共享同一批 GPU；
  \item 硬件故障、节点维护、驱动升级、镜像升级、网络异常和 checkpoint 恢复。
\end{itemize}

这时，一个核心问题出现了：谁来决定哪个任务运行在哪里？谁来保证一个 256 GPU 的训练作业不会只启动了 200 个 worker 就陷入死锁？谁来保证在线推理服务不会被离线训练任务挤占 GPU？谁来记录 GPU 利用率、HBM 占用、网络吞吐、队列等待时间和每 token 成本？谁来在节点故障后重启任务并恢复 checkpoint？

这就是 AI 集群调度系统的任务。

Kubernetes 已经成为云原生基础设施的事实标准之一。它擅长管理容器、Pod、服务发现、滚动升级、健康检查、资源声明和控制器模式。因此，很多 AI 平台选择以 Kubernetes 作为统一资源管理底座，再结合 NVIDIA Device Plugin、GPU Operator、Volcano、Kueue、Ray Operator、Kubeflow Training Operator、MPI Operator、模型服务框架和自研调度器，构建面向训练与推理的一体化 AI 集群平台。

本章不试图把 Kubernetes 所有概念完整讲一遍，而是聚焦 AI Infra 工程中最关键的调度问题：

\begin{itemize}
  \item 为什么 AI 集群不能只用普通 Kubernetes 调度策略？
  \item GPU 如何在 Kubernetes 中被发现、上报、分配和隔离？
  \item MIG、GPU sharing、拓扑感知调度分别解决什么问题？
  \item 分布式训练为什么需要 gang scheduling？
  \item Volcano 和 Kueue 在 AI batch 调度中分别扮演什么角色？
  \item LLM 推理服务如何在 Kubernetes 上部署、扩缩容和混部？
  \item AI 集群应该监控哪些指标，如何从利用率走向成本和 SLA？
\end{itemize}

\section{为什么 AI 集群需要调度系统}
\label{sec:why-ai-cluster-scheduler}

调度系统的本质是把有限资源分配给无限需求。普通在线服务调度关注 CPU、内存、端口、亲和性、反亲和性和服务可用性；AI 集群调度则更加复杂，因为 GPU 是昂贵、异构、拓扑敏感且经常以“成组”方式被使用的资源。

\subsection{GPU 资源昂贵}
\label{subsec:gpu-expensive}

GPU 是 AI 集群中最昂贵的资源之一。一台 8 GPU 服务器的采购成本、机架空间、供电、散热、网络端口和运维成本都很高。对于大模型公司或 AI 平台团队来说，GPU 空闲就是直接的资本浪费；GPU 被低效任务占用，则意味着更隐蔽的机会成本。

因此，AI 集群调度系统不能只回答“有没有空闲 GPU”，还要回答：

\begin{itemize}
  \item 空闲 GPU 是否属于同一节点、同一 NVSwitch 域或同一网络拓扑区域？
  \item 这些 GPU 的型号、显存容量、驱动版本、MIG 配置是否符合任务要求？
  \item 这个任务是在线推理还是离线训练？是否可以被抢占？
  \item 这个任务需要连续运行多久？是否支持 checkpoint 恢复？
  \item 把它放在这里会不会影响更高优先级的任务？
  \item 当前集群是否应该等待更多 GPU 空出来再启动大作业，而不是立即启动小作业造成碎片？
\end{itemize}

对于普通 Kubernetes scheduler 来说，一个 Pod 请求 1 个 GPU，调度器只要找一个有可分配 GPU 的节点即可。但对于 AI 训练，一个作业可能请求 8、16、64、256 张 GPU，并且这些 GPU 之间的拓扑关系会显著影响性能。如果只做局部贪心分配，很容易产生 GPU 碎片：表面上集群还有很多空闲 GPU，但没有足够连续、同构、拓扑良好的 GPU group 来启动大作业。

\subsection{多租户}
\label{subsec:multi-tenancy}

AI 集群通常被多个团队共享。例如基础模型团队、算法研究团队、推理平台团队、数据团队、评测团队和业务团队都需要 GPU。不同团队之间存在不同优先级、预算、配额和 SLA。

多租户调度需要解决几个问题：

\begin{itemize}
  \item \textbf{配额}：每个团队最多能使用多少 GPU、多少 GPU 小时、多少 HBM 或多少节点？
  \item \textbf{公平性}：当资源不足时，如何避免某个团队长期霸占集群？
  \item \textbf{优先级}：紧急线上修复、重要训练作业和普通实验如何排序？
  \item \textbf{抢占}：低优先级任务是否可以被高优先级任务抢占？
  \item \textbf{隔离}：不同租户之间的网络、存储、镜像、权限和日志如何隔离？
  \item \textbf{审计}：谁在什么时间使用了多少 GPU，产出了什么成本？
\end{itemize}

多租户不是简单地给 namespace 设置 resource quota。AI 任务经常是 batch job，有排队、等待、成组启动、运行数小时或数天、checkpoint 恢复等特征。调度系统需要理解“作业”而不仅仅是“Pod”。

\begin{figure}[H]
  \centering
  \begin{tikzpicture}[
    font=\small,
    box/.style={draw, rounded corners, minimum width=2.4cm, minimum height=0.75cm, align=center},
    pool/.style={draw, rounded corners, minimum width=9.8cm, minimum height=2.2cm, align=center},
    arrow/.style={-{Stealth[length=2mm]}, thick}
  ]
    \node[box] (teamA) at (-4,4.0) {基础模型团队};
    \node[box] (teamB) at (-1.3,4.0) {推理平台团队};
    \node[box] (teamC) at (1.4,4.0) {评测团队};
    \node[box] (teamD) at (4.1,4.0) {数据团队};

    \node[box, minimum width=8.8cm] (queue) at (0,2.7) {队列 / 配额 / 优先级 / 抢占策略};

    \node[pool] (cluster) at (0,1.0) {};
    \node at (0,1.75) {AI GPU 资源池};
    \foreach \i/\x in {0/-4.0,1/-3.0,2/-2.0,3/-1.0,4/0.0,5/1.0,6/2.0,7/3.0,8/4.0} {
      \node[draw, rounded corners, minimum width=0.65cm, minimum height=0.45cm] at (\x,0.65) {G};
    }

    \draw[arrow] (teamA.south) -- (queue.north);
    \draw[arrow] (teamB.south) -- (queue.north);
    \draw[arrow] (teamC.south) -- (queue.north);
    \draw[arrow] (teamD.south) -- (queue.north);
    \draw[arrow] (queue.south) -- (cluster.north);

    \node[align=center] at (0,-0.55) {多租户 AI 集群调度：不同团队通过队列、配额和优先级共享 GPU 资源池。};
  \end{tikzpicture}
  \caption{多租户 AI 集群资源池：调度系统需要同时考虑配额、公平性、优先级和资源碎片。}
  \label{fig:multi-tenant-ai-cluster}
\end{figure}

\subsection{训练、推理、评测混部}
\label{subsec:training-serving-eval-colocation}

AI 集群中的任务类型差异很大。预训练任务通常是长时间、大规模、通信密集型；推理服务是长期运行、延迟敏感、需要弹性伸缩；评测任务可能是短时突发、吞吐型；数据处理可能主要消耗 CPU、内存、存储和网络；微调任务规模较小但数量很多。

不同任务对调度系统的要求不同：

\begin{table}[H]
  \centering
  \caption{AI 集群中不同任务类型的调度特征}
  \label{tab:ai-workload-scheduling-characteristics}
  \begin{tabular}{p{2.8cm}p{3.5cm}p{4.2cm}p{4.0cm}}
    \toprule
    任务类型 & 资源特征 & 调度重点 & 典型风险 \\
    \midrule
    预训练 & 大规模 GPU、长时间、强通信 & gang scheduling、拓扑感知、checkpoint 恢复 & 启动碎片、通信慢、故障重启成本高 \\
    SFT / 微调 & 中小规模 GPU、任务数量多 & 队列、公平性、镜像和数据准备 & 小任务挤占大作业资源、低利用率 \\
    RLHF / DPO & 多阶段、多角色、可能混合训练和推理 & workflow 编排、actor/rollout/reference 模型资源协调 & 阶段不均衡、资源等待 \\
    在线推理 & 长期运行、延迟敏感、显存常驻 & 高可用、扩缩容、bin packing、SLA & 冷启动、尾延迟、显存碎片 \\
    离线评测 & 短时突发、吞吐型 & 快速启动、低优先级抢占、批量排队 & 影响线上服务或重要训练 \\
    数据处理 & CPU、内存、存储、网络为主 & 与 GPU 任务隔离、I/O 调度 & 抢占网络和存储带宽 \\
    \bottomrule
  \end{tabular}
\end{table}

混部可以提高资源利用率，但也会增加干扰。例如在线推理服务可能 SM utilization 不高，但 HBM 占用高且对延迟敏感；训练作业可能吞吐高但通信突发；评测任务可能频繁启动和释放 GPU，造成资源碎片。调度系统需要根据任务特征做隔离和分层，而不是把所有 GPU workload 当成同一种 Pod。

\subsection{利用率与公平性}
\label{subsec:utilization-and-fairness}

AI 集群调度的两个目标经常冲突：提高利用率和保证公平性。为了提高利用率，调度器希望尽快把空闲 GPU 分配出去；为了公平性和大作业启动，调度器又可能需要让一些小任务等待，以便积累足够连续资源给大任务。

例如，集群当前有 4 个节点，每个节点剩余 2 张 GPU，总共 8 张空闲 GPU。一个 8 GPU 单节点 TP 推理实例无法启动，因为它需要同一节点的 8 张 GPU；一个 8 GPU 训练作业如果要求 NVSwitch 单节点，也无法启动。此时从总量看资源足够，从拓扑和连续性看资源不足。这就是 GPU 碎片。

调度系统需要在以下策略之间权衡：

\begin{itemize}
  \item \textbf{bin packing}：尽量把任务塞满少数节点，减少碎片；
  \item \textbf{spread}：把任务分散到多个节点，提高可用性或降低局部干扰；
  \item \textbf{reservation}：为大作业预留资源，避免小作业打散集群；
  \item \textbf{backfill}：在不影响高优先级大作业预计启动时间的前提下，运行短小任务；
  \item \textbf{preemption}：通过抢占低优先级任务释放资源；
  \item \textbf{quota borrowing}：允许队列临时借用其他队列未使用的配额。
\end{itemize}

优秀的 AI 调度系统不是简单地追求某一时刻 GPU utilization 最高，而是要优化更长期的目标：整体吞吐、队列等待时间、SLA、训练完成时间、公平性和成本。

\section{Kubernetes 基础回顾}
\label{sec:kubernetes-basics-review}

本节快速回顾 Kubernetes 中与 AI 调度关系最密切的几个概念：Pod、Node、Scheduler、Device Plugin 和 Operator。熟悉 Kubernetes 的读者可以快速浏览。

\subsection{Pod}
\label{subsec:pod}

Pod 是 Kubernetes 中最小的调度单元。一个 Pod 可以包含一个或多个容器，这些容器共享网络 namespace、部分存储卷和生命周期。在 AI 任务中，一个 Pod 可能代表：

\begin{itemize}
  \item 一个训练 worker；
  \item 一个 parameter server；
  \item 一个 PyTorch distributed rank；
  \item 一个 Ray worker；
  \item 一个 vLLM / TensorRT-LLM / Triton model server 实例；
  \item 一个数据预处理任务；
  \item 一个评测 worker。
\end{itemize}

Pod 通过 \texttt{resources.requests} 和 \texttt{resources.limits} 声明资源需求。GPU 通常以扩展资源的形式出现，例如 \texttt{nvidia.com/gpu: 1}。一个最小 GPU Pod 示例：

\begin{lstlisting}[caption={请求 1 张 GPU 的 Pod 示例},label={lst:k8s-gpu-pod}]
apiVersion: v1
kind: Pod
metadata:
  name: gpu-demo
spec:
  restartPolicy: Never
  containers:
  - name: app
    image: nvcr.io/nvidia/pytorch:xx.xx-py3
    command: ["python", "train.py"]
    resources:
      limits:
        nvidia.com/gpu: 1
\end{lstlisting}

在 Kubernetes 中，GPU 这类设备资源通常要求 request 与 limit 相等，常见写法是只写 limits。实际行为取决于 device plugin 和 Kubernetes 版本，但工程上应避免对 GPU 做类似 CPU 那样的 overcommit，除非明确使用 GPU sharing 或 MIG 等机制。

\subsection{Node}
\label{subsec:node}

Node 是 Kubernetes 集群中的工作节点。对 AI 集群来说，一个 Node 通常是一台 GPU 服务器，包含：

\begin{itemize}
  \item 多张 GPU；
  \item CPU 和内存；
  \item 本地 NVMe；
  \item 网卡和 RDMA 设备；
  \item NVIDIA driver、CUDA runtime、container runtime；
  \item kubelet、device plugin、monitoring agent；
  \item 可能还包括 MIG 配置、拓扑标签和健康检查组件。
\end{itemize}

Node 会向 Kubernetes API Server 上报可分配资源。调度器根据 Pod 的资源请求、node selector、affinity、taint/toleration、topology spread、priority 等信息选择目标节点。

AI 集群中常见的 node label 包括：

\begin{lstlisting}[language=bash,caption={AI 集群常见 Node Label 示例},label={lst:ai-node-labels}]
accelerator=nvidia
gpu.vendor=nvidia
gpu.model=h100
gpu.count=8
gpu.memory=80gb
gpu.fabric=nvswitch
rdma=true
network.rail=4
zone=az-a
pool=training
\end{lstlisting}

这些 label 可以帮助调度器或平台层把任务放到合适资源池。例如，预训练任务只允许调度到 \texttt{pool=training} 且 \texttt{gpu.fabric=nvswitch} 的节点；在线推理任务调度到 \texttt{pool=serving}；评测任务可以使用较低优先级节点池。

\subsection{Scheduler}
\label{subsec:scheduler}

Kubernetes Scheduler 的基本职责是为 Pending Pod 选择一个 Node。它通常经历两个阶段：

\begin{enumerate}
  \item \textbf{Filter}：过滤掉不满足资源、亲和性、污点容忍、端口、卷等条件的节点；
  \item \textbf{Score}：对可行节点打分，选择最优节点。
\end{enumerate}

默认 scheduler 对通用 workload 很有效，但对 AI batch workload 有几个不足：

\begin{itemize}
  \item 它以 Pod 为基本单位，不天然理解一个分布式训练 Job 需要所有 worker 一起启动；
  \item 它不一定理解 GPU 拓扑、NVLink、NIC 亲和性、MIG profile 等细节；
  \item 它不提供复杂队列、配额借用、gang scheduling 和 backfill 能力；
  \item 它对长时间运行的 batch job 与在线 service 的混部策略有限；
  \item 它不直接处理 checkpoint-aware preemption 和作业级恢复。
\end{itemize}

因此，AI 集群常常在默认 scheduler 之外引入 Volcano、Kueue、自研 scheduler extender、scheduler plugin 或更上层的 batch queueing system。

\subsection{Device Plugin}
\label{subsec:device-plugin}

Device Plugin 是 Kubernetes 管理 GPU、NIC、FPGA 等特殊硬件资源的扩展机制。设备厂商或平台团队通过 device plugin 向 kubelet 注册资源类型、上报设备列表，并在 Pod 被分配设备时执行必要的环境变量、设备文件、挂载和 runtime 配置。

以 NVIDIA GPU 为例，NVIDIA device plugin 会向 Kubernetes 暴露类似 \texttt{nvidia.com/gpu} 的资源。Pod 请求该资源后，kubelet 会调用 device plugin 分配具体 GPU，并把对应设备注入容器。

Device Plugin 的重要性在于：Kubernetes 本身并不知道 GPU 内部细节。它需要 device plugin 告诉它节点上有哪些设备、哪些设备健康、如何把设备分配给容器。对于 AI Infra 工程师来说，device plugin 是 GPU 从“物理硬件”变成“Kubernetes 可调度资源”的关键桥梁。

\subsection{Operator}
\label{subsec:operator}

Operator 是 Kubernetes 的一种控制器模式，用自定义资源 CRD 和控制循环管理复杂应用。AI 集群中常见 operator 包括：

\begin{itemize}
  \item NVIDIA GPU Operator：安装和管理 driver、device plugin、DCGM exporter、MIG manager 等 GPU 组件；
  \item Training Operator / Kubeflow Operator：管理 PyTorchJob、TFJob、MPIJob 等分布式训练任务；
  \item Ray Operator：管理 Ray cluster 和 Ray job；
  \item Model serving operator：管理模型部署、版本、扩缩容和流量切换；
  \item Volcano / Kueue 相关控制器：管理队列、作业准入和调度。
\end{itemize}

Operator 的价值是把复杂运维动作自动化。例如 GPU 节点加入集群后，GPU Operator 可以自动部署驱动相关组件和监控 exporter；用户提交一个 PyTorchJob 后，训练 operator 可以创建 master 和 worker Pod，配置 rendezvous，监控状态，并在失败时重试。

\begin{figure}[H]
  \centering
  \begin{tikzpicture}[
    font=\small,
    box/.style={draw, rounded corners, minimum width=2.2cm, minimum height=0.75cm, align=center},
    wide/.style={draw, rounded corners, minimum width=8.8cm, minimum height=0.85cm, align=center},
    nodebox/.style={draw, rounded corners, minimum width=8.8cm, minimum height=2.1cm, align=center},
    arrow/.style={-{Stealth[length=2mm]}, thick}
  ]
    \node[wide] (api) at (0,4.7) {API Server};
    \node[box] (sched) at (-3.2,3.4) {Scheduler};
    \node[box] (ctrl) at (0,3.4) {Controller\\Manager};
    \node[box] (op) at (3.2,3.4) {Operators};

    \node[nodebox] (node) at (0,1.55) {};
    \node[anchor=north west] at ($(node.north west)+(0.2,-0.2)$) {GPU Node};
    \node[box] (kubelet) at (-3.0,1.4) {kubelet};
    \node[box] (dp) at (0,1.4) {Device\\Plugin};
    \node[box] (pod) at (3.0,1.4) {GPU Pods};

    \node[box, minimum width=7.5cm] (runtime) at (0,0.35) {Container Runtime / NVIDIA Runtime / Driver};

    \draw[arrow] (sched) -- (api);
    \draw[arrow] (ctrl) -- (api);
    \draw[arrow] (op) -- (api);
    \draw[arrow] (api) -- (kubelet);
    \draw[arrow] (dp) -- (kubelet);
    \draw[arrow] (kubelet) -- (pod);
    \draw[arrow] (pod) -- (runtime);

    \node[align=center] at (0,-0.65) {Kubernetes 控制面通过 kubelet、device plugin 和 operator 管理 GPU 节点与 AI workload。};
  \end{tikzpicture}
  \caption{Kubernetes 控制面与 GPU 节点组件关系图。}
  \label{fig:k8s-control-plane-gpu-node}
\end{figure}

\section{GPU 资源管理}
\label{sec:gpu-resource-management}

GPU 资源管理是 Kubernetes AI 平台的核心。它不仅包括“让 Pod 能请求 GPU”，还包括设备健康检查、驱动管理、MIG、GPU sharing、拓扑感知、资源隔离、监控和故障处理。

\subsection{NVIDIA Device Plugin}
\label{subsec:nvidia-device-plugin}

NVIDIA Device Plugin 的主要职责是把节点上的 NVIDIA GPU 暴露给 Kubernetes。它通常以 DaemonSet 形式部署在 GPU 节点上，与 kubelet 通信。工作流程可以简化为：

\begin{enumerate}
  \item device plugin 启动后发现节点上的 GPU；
  \item 向 kubelet 注册资源名，例如 \texttt{nvidia.com/gpu}；
  \item 周期性上报设备健康状态；
  \item 当 Pod 请求 GPU 时，kubelet 调用 device plugin 的 Allocate 接口；
  \item device plugin 返回容器需要的设备文件、环境变量和挂载信息；
  \item 容器启动后只能看到被分配的 GPU。
\end{enumerate}

从平台角度看，需要关注以下问题：

\begin{itemize}
  \item device plugin 版本是否与驱动、container runtime、Kubernetes 版本兼容；
  \item GPU 健康状态是否能正确反映到 Kubernetes；
  \item GPU reset、Xid error、ECC error 后是否自动隔离节点；
  \item 容器中 \texttt{CUDA\_VISIBLE\_DEVICES} 是否符合预期；
  \item MIG 模式下资源名和 profile 是否正确上报；
  \item 与 GPU Operator、DCGM exporter、MIG manager 是否协同。
\end{itemize}

\subsection{GPU Request / Limit}
\label{subsec:gpu-request-limit}

用户通常通过 Pod spec 声明 GPU 需求：

\begin{lstlisting}[caption={在 Kubernetes 中请求 GPU 资源},label={lst:k8s-gpu-resource-request}]
resources:
  limits:
    nvidia.com/gpu: 4
\end{lstlisting}

这里的 \texttt{4} 表示该容器需要 4 张 GPU。Kubernetes 会把该 Pod 调度到一个可分配至少 4 张 GPU 的节点上。对于分布式训练，通常每个 Pod 请求 1 张或多张 GPU，然后由训练框架根据 rank、world size 和 local rank 建立通信。

一个 PyTorch 分布式训练 worker 的 Pod 可能类似：

\begin{lstlisting}[caption={分布式训练 Worker Pod 的资源声明示例},label={lst:pytorch-worker-gpu-request}]
apiVersion: v1
kind: Pod
metadata:
  name: train-worker-0
spec:
  restartPolicy: Never
  containers:
  - name: worker
    image: my-registry/pytorch-train:latest
    command: ["bash", "-lc", "torchrun --nproc_per_node=8 train.py"]
    resources:
      limits:
        nvidia.com/gpu: 8
        cpu: "64"
        memory: "512Gi"
\end{lstlisting}

需要注意，Kubernetes 原生资源声明并不表达 GPU 拓扑。例如请求 8 张 GPU，只表示需要同一节点上 8 个 GPU 设备，并不说明这些 GPU 是否通过 NVSwitch 互联，也不说明 GPU 与 NIC 的关系。要表达这些约束，需要结合 node label、node affinity、scheduler plugin、device plugin 拓扑信息或自研调度层。

\subsection{MIG}
\label{subsec:mig}

MIG，即 Multi-Instance GPU，是 NVIDIA 在部分数据中心 GPU 上提供的硬件级 GPU 切分能力。它可以把一张物理 GPU 切分成多个 GPU instance，每个 instance 拥有独立的一部分计算资源、显存、缓存和带宽隔离能力。MIG 的目标不是提高单个大训练任务性能，而是提高小任务、小模型推理和多租户场景下的资源利用率与隔离性。

例如，一张大显存 GPU 如果直接分配给一个只需要少量显存和算力的小模型，会造成浪费。通过 MIG，可以把它切分成多个较小 profile，让多个推理实例或轻量任务共享同一张物理 GPU，同时获得比普通进程共享更强的硬件隔离。

\begin{figure}[H]
  \centering
  \begin{tikzpicture}[
    font=\small,
    gpu/.style={draw, rounded corners, minimum width=7.8cm, minimum height=2.5cm, align=center},
    slice/.style={draw, rounded corners, minimum width=1.2cm, minimum height=1.7cm, align=center, fill=gray!10},
    arrow/.style={-{Stealth[length=2mm]}, thick}
  ]
    \node[gpu] (full) at (0,3.2) {};
    \node at (0,4.15) {Physical GPU};
    \node at (0,3.2) {完整 GPU：SM / HBM / L2 / Copy Engine};

    \node[gpu] (mig) at (0,0.7) {};
    \node at (0,1.65) {MIG Mode};
    \foreach \i/\x/\txt in {
      0/-3.0/{GI 0},
      1/-1.8/{GI 1},
      2/-0.6/{GI 2},
      3/0.6/{GI 3},
      4/1.8/{GI 4},
      5/3.0/{GI 5}
    } {
      \node[slice] at (\x,0.65) {\txt\\SM\\HBM};
    }

    \draw[arrow] (full.south) -- (mig.north);
    \node[align=center] at (0,-1.0) {MIG 将一张物理 GPU 切分为多个具有硬件隔离的 GPU instance。};
  \end{tikzpicture}
  \caption{MIG GPU 切分示意图：适合小模型推理、多租户隔离和提高碎片资源利用率。}
  \label{fig:mig-gpu-partition}
\end{figure}

MIG 在 Kubernetes 中通常由 GPU Operator 和 MIG Manager 管理。平台可以为节点设置不同 MIG strategy，例如单一 profile 或混合 profile。调度时，MIG instance 会以不同扩展资源的形式暴露给 Kubernetes，例如某种 \texttt{nvidia.com/mig-...} 资源名。

MIG 的使用需要权衡：

\begin{itemize}
  \item \textbf{优点}：硬件隔离更强，适合多租户小模型推理，提高大 GPU 的切分利用率；
  \item \textbf{缺点}：切分后无法灵活合并，profile 选择不当会造成新的碎片；
  \item \textbf{限制}：并非所有 GPU 支持，且某些多卡通信、NVLink、profiling 或共享场景会有额外约束；
  \item \textbf{运维复杂度}：需要管理 MIG 模式切换、profile 配置、节点 drain、资源上报和监控。
\end{itemize}

实践中，MIG 更适合稳定的推理资源池，而不是频繁变化的大规模训练资源池。对于训练，通常更希望保留完整 GPU，以获得完整算力、显存和互联能力。

\subsection{GPU Sharing}
\label{subsec:gpu-sharing}

GPU sharing 指多个工作负载共享同一张 GPU。它与 MIG 不同，MIG 是硬件切分，GPU sharing 更多是软件或 runtime 层的共享。常见方式包括：

\begin{itemize}
  \item 多个进程直接共享同一 GPU；
  \item NVIDIA MPS；
  \item time-slicing；
  \item framework-level batching；
  \item serving runtime 内部的多模型复用；
  \item MIG + sharing 的组合。
\end{itemize}

GPU sharing 的价值在于提高小任务利用率。例如某些 notebook、开发调试、小模型推理、embedding 服务或低 QPS 模型不需要独占整张 GPU。共享可以减少浪费。

但 GPU sharing 风险也很大：

\begin{itemize}
  \item 显存 OOM 会相互影响；
  \item SM、HBM 带宽、L2 cache、copy engine 竞争导致性能抖动；
  \item 在线推理尾延迟变差；
  \item 监控和计费复杂；
  \item 故障隔离弱于 MIG；
  \item 用户容易误以为自己独占 GPU。
\end{itemize}

因此，GPU sharing 更适合低优先级、开发测试或可容忍抖动的推理场景。对于强 SLA 在线推理和大规模训练，应谨慎使用。

\subsection{Topology-Aware Scheduling}
\label{subsec:topology-aware-scheduling}

拓扑感知调度是 AI 集群调度的关键能力。它需要理解 GPU 与 GPU、GPU 与 NIC、GPU 与 CPU NUMA、节点与交换机之间的关系。前两章已经解释过，TP、EP、RDMA 和 GPUDirect RDMA 都强依赖拓扑。

在 Kubernetes 中，可以通过多种方式实现拓扑感知：

\begin{itemize}
  \item 节点标签：例如 \texttt{gpu.fabric=nvswitch}、\texttt{rdma=true}；
  \item node affinity：限制 Pod 只能调度到满足拓扑条件的节点；
  \item pod affinity / anti-affinity：控制 Pod 之间靠近或分散；
  \item topology spread constraints：控制副本在 zone、rack、node 间分布；
  \item device plugin 上报拓扑信息；
  \item scheduler extender 或 scheduler plugin；
  \item 上层 batch scheduler 进行作业级节点选择；
  \item 自研资源画像系统记录 GPU/NIC/NVLink/NVSwitch 拓扑。
\end{itemize}

\begin{figure}[H]
  \centering
  \begin{tikzpicture}[
    font=\small,
    gpu/.style={draw, rounded corners, minimum width=0.9cm, minimum height=0.5cm, align=center},
    group/.style={draw, rounded corners, dashed, inner sep=0.22cm},
    task/.style={draw, rounded corners, minimum width=2.0cm, minimum height=0.65cm, align=center, fill=gray!10},
    link/.style={line width=0.8pt},
    arrow/.style={-{Stealth[length=2mm]}, thick}
  ]
    \foreach \i/\x in {0/-3.5,1/-2.5,2/-1.5,3/-0.5,4/0.7,5/1.7,6/2.7,7/3.7} {
      \node[gpu] (g\i) at (\x,0.6) {G\i};
    }

    \node[group, fit=(g0)(g1)(g2)(g3), label=below:{NVSwitch Group A}] {};
    \node[group, fit=(g4)(g5)(g6)(g7), label=below:{NVSwitch Group B}] {};

    \node[task] (job) at (0,3.1) {TP=4 推理实例};

    \draw[arrow] (job.south) -- (-2.0,1.25);
    \node[align=center] at (0,2.25) {拓扑感知调度选择同一高速互联 group 内的 GPU};

    \foreach \a/\b in {0/1,1/2,2/3,4/5,5/6,6/7} {
      \draw[link] (g\a) -- (g\b);
    }
  \end{tikzpicture}
  \caption{拓扑感知 GPU 分配示意：通信密集任务应优先绑定同一高速互联域内的 GPU。}
  \label{fig:topology-aware-gpu-allocation}
\end{figure}

拓扑感知调度尤其适用于：

\begin{itemize}
  \item TP 推理实例；
  \item 单节点 8 GPU 训练；
  \item 多节点训练中的 GPU-NIC 亲和；
  \item MoE expert parallel；
  \item 多 rail RDMA；
  \item 同机多模型混部。
\end{itemize}

\section{分布式训练作业调度}
\label{sec:distributed-training-scheduling}

分布式训练与普通 Pod 调度最大的区别是：它通常需要一组 Pod 共同启动、共同运行、共同失败或共同成功。一个 64 GPU PyTorch 训练作业如果只启动了 50 个 worker，剩下 14 个 worker 处于 Pending，已经启动的 worker 可能会一直等待 rendezvous，浪费 GPU。这个问题就是 gang scheduling 的动机。

\subsection{Gang Scheduling}
\label{subsec:gang-scheduling}

Gang scheduling，也称 coscheduling，要求一个作业的一组任务要么一起调度成功，要么都不运行。对于分布式训练来说，这非常重要。一个 PyTorchJob、MPIJob 或 Ray cluster 往往有最小可运行副本数。如果资源不够，调度器应该让整个 job 等待，而不是只启动一部分 Pod。

\begin{figure}[H]
  \centering
  \begin{tikzpicture}[
    font=\small,
    pod/.style={draw, rounded corners, minimum width=0.85cm, minimum height=0.45cm, align=center},
    node/.style={draw, rounded corners, minimum width=2.0cm, minimum height=0.7cm, align=center},
    arrow/.style={-{Stealth[length=2mm]}, thick}
  ]
    \node at (0,4.0) {Job 需要 8 个 worker，minAvailable = 8};

    \foreach \i/\x in {0/-3.5,1/-2.5,2/-1.5,3/-0.5,4/0.5,5/1.5,6/2.5,7/3.5} {
      \node[pod] (p\i) at (\x,3.2) {W\i};
    }

    \node[node] (n0) at (-3,1.7) {Node A\\4 GPU};
    \node[node] (n1) at (-0.5,1.7) {Node B\\2 GPU};
    \node[node] (n2) at (2.0,1.7) {Node C\\1 GPU};

    \draw[arrow, dashed] (p0.south) -- (n0.north);
    \draw[arrow, dashed] (p1.south) -- (n0.north);
    \draw[arrow, dashed] (p2.south) -- (n0.north);
    \draw[arrow, dashed] (p3.south) -- (n0.north);
    \draw[arrow, dashed] (p4.south) -- (n1.north);
    \draw[arrow, dashed] (p5.south) -- (n1.north);
    \draw[arrow, dashed] (p6.south) -- (n2.north);

    \node[draw, rounded corners, minimum width=8.5cm, minimum height=0.7cm, align=center, fill=gray!10] at (0,0.35) {资源不足 8 个 worker：整个 Job 等待，避免部分启动浪费 GPU};
  \end{tikzpicture}
  \caption{Gang scheduling 等待与启动流程：资源不足时不应只启动部分训练 worker。}
  \label{fig:gang-scheduling}
\end{figure}

Gang scheduling 需要调度器理解 job 级语义，而不是只看单个 Pod。它通常与队列、配额、优先级和抢占结合使用。例如一个高优先级 128 GPU 训练 job 等待时，调度器可能需要抢占一些低优先级任务，释放足够资源后一次性启动整个 job。

\subsection{Volcano}
\label{subsec:volcano}

Volcano 是面向云原生 batch 和高性能计算 workload 的调度系统。它与 Kubernetes 兼容，并提供队列、gang scheduling、优先级、抢占、资源回收、backfill、DRF、公平调度等能力。对 AI 训练来说，Volcano 常用于调度 PyTorch、TensorFlow、MPI、Spark、Ray 等计算密集型任务。

Volcano 的核心价值在于它把调度对象从单个 Pod 提升到 Job 和 Queue。一个 VolcanoJob 可以包含多个 task，每个 task 可以有副本数、资源需求和最小可用数量。调度器根据队列和作业约束决定是否启动。

一个简化 VolcanoJob 示例：

\begin{lstlisting}[caption={VolcanoJob 的概念示例},label={lst:volcano-job-example}]
apiVersion: batch.volcano.sh/v1alpha1
kind: Job
metadata:
  name: pretrain-job
spec:
  minAvailable: 8
  schedulerName: volcano
  queue: foundation-model
  tasks:
  - name: worker
    replicas: 8
    template:
      spec:
        restartPolicy: Never
        containers:
        - name: trainer
          image: my-registry/pretrain:latest
          command: ["bash", "-lc", "torchrun --nnodes=1 --nproc_per_node=8 train.py"]
          resources:
            limits:
              nvidia.com/gpu: 1
\end{lstlisting}

这个示例表达了一个重要语义：作业至少需要 8 个 worker 才能运行。如果只调度了 3 个 worker，它没有意义。Volcano 的 gang plugin 可以帮助避免这种部分启动问题。

Volcano 适合的场景包括：

\begin{itemize}
  \item 大规模分布式训练；
  \item 多队列、多租户 batch 集群；
  \item 需要 gang scheduling 的 MPI/PyTorch/Ray workload；
  \item 需要公平调度、优先级、抢占和 backfill 的场景；
  \item HPC 与 AI 混合计算集群。
\end{itemize}

\subsection{Kueue}
\label{subsec:kueue}

Kueue 是 Kubernetes 原生的作业排队系统，面向 batch、HPC、AI/ML 等 workload。它的核心思想不是替代所有调度逻辑，而是在 Pod 真正创建或运行之前，对作业进行准入控制：根据队列、配额、资源可用性和优先级判断作业什么时候可以进入集群运行。

Kueue 中常见概念包括：

\begin{itemize}
  \item \textbf{Workload}：代表一个待调度的作业抽象；
  \item \textbf{LocalQueue}：namespace 内用户提交作业的队列；
  \item \textbf{ClusterQueue}：跨 namespace 的资源配额和队列管理对象；
  \item \textbf{ResourceFlavor}：不同类型资源，例如不同 GPU 型号、不同节点池；
  \item \textbf{Admission}：当资源和配额满足时，Kueue 接纳作业运行；
  \item \textbf{Preemption}：在策略允许时，抢占低优先级 workload。
\end{itemize}

Kueue 的一个典型价值是把“排队”和“运行”分开。用户可以提交大量训练或评测任务，但只有被接纳的任务才真正创建大量 Pod 消耗资源。这对于控制集群压力、实现配额、公平性和多租户非常重要。

一个简化的 Kueue 资源概念如下：

\begin{lstlisting}[caption={Kueue ClusterQueue 与 ResourceFlavor 的概念示例},label={lst:kueue-clusterqueue-example}]
apiVersion: kueue.x-k8s.io/v1beta1
kind: ResourceFlavor
metadata:
  name: h100-nvswitch
---
apiVersion: kueue.x-k8s.io/v1beta1
kind: ClusterQueue
metadata:
  name: foundation-model-queue
spec:
  namespaceSelector: {}
  resourceGroups:
  - coveredResources: ["cpu", "memory", "nvidia.com/gpu"]
    flavors:
    - name: h100-nvswitch
      resources:
      - name: nvidia.com/gpu
        nominalQuota: 256
      - name: cpu
        nominalQuota: 2048
      - name: memory
        nominalQuota: 16Ti
\end{lstlisting}

Kueue 适合与 Job、RayJob、MPIJob、PyTorchJob 等作业控制器结合，作为统一准入和配额层。它尤其适合多团队共享 batch 集群，希望按照资源 flavor 和队列控制不同团队使用资源的场景。

\subsection{Priority 与 Preemption}
\label{subsec:priority-preemption}

优先级和抢占是调度系统必须谨慎设计的能力。没有优先级，紧急任务无法快速获得资源；没有抢占，高优先级任务可能长期等待；但如果抢占过于激进，低优先级任务可能永远无法完成，集群也会频繁浪费已经完成的训练进度。

AI 任务的抢占策略需要结合 checkpoint。对于支持 checkpoint 的训练任务，可以在抢占前发送信号，让任务保存状态；对于不支持 checkpoint 的任务，抢占意味着从头开始，成本很高。对于在线推理服务，抢占必须考虑流量迁移、连接 draining 和 SLA。

常见策略包括：

\begin{itemize}
  \item 高优先级线上服务不被抢占；
  \item 低优先级实验任务可被抢占；
  \item 预训练任务只有在 checkpoint 完成后才允许安全抢占；
  \item 抢占前设置 grace period；
  \item 被抢占任务重新进入队列；
  \item 对频繁被抢占的任务提供 aging 或补偿机制；
  \item 抢占应尽量释放连续资源，而不是随机杀 Pod。
\end{itemize}

\subsection{Checkpoint-Based Rescheduling}
\label{subsec:checkpoint-based-rescheduling}

大规模训练不可能假设节点永远不坏。GPU Xid、ECC error、网络抖动、节点重启、驱动问题、存储故障、作业代码异常都可能导致训练中断。因此，调度系统必须与 checkpoint 机制结合。

Checkpoint-based rescheduling 的流程通常是：

\begin{enumerate}
  \item 训练任务周期性保存 distributed checkpoint；
  \item 调度系统监控 Pod、Node 和 Job 状态；
  \item 节点故障或抢占发生时，Job 标记为失败或待恢复；
  \item 调度器重新为 Job 分配资源；
  \item 训练任务从最近 checkpoint 恢复；
  \item 平台记录失败原因、恢复耗时和损失的训练步数。
\end{enumerate}

对大模型训练来说，checkpoint 本身也是重 I/O 操作。如果所有大作业同时 checkpoint，会冲击存储和网络。调度系统可以参与 checkpoint policy，例如错峰 checkpoint、限制并发 checkpoint、根据作业优先级调整 checkpoint 间隔。

\section{推理服务编排}
\label{sec:serving-orchestration}

训练作业通常是 batch workload，而推理服务是 online workload。它们都使用 GPU，但调度目标完全不同。训练追求吞吐和完成时间，推理追求可用性、延迟、吞吐和成本。Kubernetes 原生适合管理长期运行服务，因此很多 LLM Serving 系统会部署在 Kubernetes 上，但 GPU 推理服务仍然有大量特殊问题。

\subsection{Model Server Deployment}
\label{subsec:model-server-deployment}

一个 LLM 推理服务通常包含：

\begin{itemize}
  \item 前端 API gateway；
  \item tokenizer / request processor；
  \item model server，例如 vLLM、TensorRT-LLM、Triton、Text Generation Inference 或自研 runtime；
  \item scheduler，用于 continuous batching、prefill/decode 调度和 admission control；
  \item KV Cache manager；
  \item metrics exporter；
  \item model artifact loader；
  \item 日志、trace 和计费组件。
\end{itemize}

在 Kubernetes 中，最简单的部署方式是使用 Deployment 或 StatefulSet 管理模型服务 Pod，用 Service 暴露访问入口。对于需要多 GPU tensor parallel 的模型，一个 Pod 可能请求 2、4 或 8 张 GPU，并在容器内部启动多进程 runtime。

\begin{lstlisting}[caption={LLM Model Server Deployment 的概念示例},label={lst:llm-serving-deployment}]
apiVersion: apps/v1
kind: Deployment
metadata:
  name: llama-serving
spec:
  replicas: 2
  selector:
    matchLabels:
      app: llama-serving
  template:
    metadata:
      labels:
        app: llama-serving
    spec:
      nodeSelector:
        pool: serving
        gpu.fabric: nvswitch
      containers:
      - name: server
        image: my-registry/llm-serving:latest
        command: ["bash", "-lc", "python -m server --model /models/llama --tensor-parallel-size 4"]
        ports:
        - containerPort: 8000
        resources:
          limits:
            nvidia.com/gpu: 4
            cpu: "32"
            memory: "256Gi"
\end{lstlisting}

需要注意，Deployment 的滚动升级语义对大模型服务未必总是合适。模型加载可能需要数分钟，显存初始化和 CUDA graph capture 也需要时间。如果滚动策略设置不当，可能造成可用副本不足或冷启动抖动。对于大模型服务，通常需要自定义 readiness probe、预热流程、流量切换和灰度发布策略。

\subsection{Autoscaling}
\label{subsec:autoscaling}

普通 Web 服务常用 CPU utilization 或 QPS 做 autoscaling，但 LLM 推理服务的扩缩容指标更复杂。常见指标包括：

\begin{itemize}
  \item QPS；
  \item active requests；
  \item waiting queue length；
  \item TTFT；
  \item TPOT；
  \item request latency P50/P95/P99；
  \item GPU utilization；
  \item HBM usage；
  \item KV Cache usage；
  \item tokens per second；
  \item prefill queue length 与 decode queue length；
  \item SLA violation rate。
\end{itemize}

LLM serving 的扩缩容难点在于冷启动成本高。一个大模型实例启动需要拉取镜像、加载权重、初始化 CUDA context、构建通信组、预热 kernel，甚至加载量化权重和构建 CUDA graph。这个过程可能远慢于普通服务。因此，autoscaling 不能只在流量升高后才启动新实例，还需要预测、预留和 warm pool。

推理扩缩容的常见策略包括：

\begin{itemize}
  \item 保留最小 warm replicas；
  \item 根据 queue length 和 token throughput 扩容；
  \item 对 prefill-heavy 和 decode-heavy 资源池分别扩缩容；
  \item 使用 admission control 防止队列无限增长；
  \item 在流量低谷缩容，但保留热点模型；
  \item 对多模型服务做模型级别的容量规划。
\end{itemize}

\subsection{GPU Bin Packing}
\label{subsec:gpu-bin-packing}

推理服务通常比训练更容易产生 GPU 碎片。原因包括：

\begin{itemize}
  \item 不同模型大小不同，显存需求不同；
  \item 不同模型 QPS 不同，SM 和 HBM 带宽使用率不同；
  \item 一些模型需要 1 张 GPU，一些需要 4 张或 8 张 GPU；
  \item LoRA adapter、多租户模型和 embedding/rerank 服务共存；
  \item 在线服务不能随意迁移，否则会影响 SLA。
\end{itemize}

Bin packing 的目标是把模型实例合理放置到 GPU 节点上，提高资源利用率并减少碎片。例如，一个 8 GPU 节点可以放置两个 TP=4 的模型实例，也可以放置一个 TP=8 的大模型实例，还可以放置多个 MIG 小实例。不同放置方式对显存、通信、带宽和可用性影响不同。

推理 bin packing 需要考虑：

\begin{itemize}
  \item GPU 显存；
  \item HBM 带宽；
  \item SM utilization；
  \item KV Cache 峰值；
  \item 模型权重常驻；
  \item TP group 拓扑；
  \item 实例之间的干扰；
  \item 负载均衡和故障域。
\end{itemize}

如果只按显存打包，可能把多个 memory-bound decode 服务放在同一节点，导致 HBM 带宽争用和尾延迟上升。如果只按 SM utilization 打包，也可能忽略 KV Cache 和网络通信。因此，高质量推理调度需要 workload profiling。

\subsection{Cold Start}
\label{subsec:cold-start}

大模型推理的冷启动是一个严重问题。普通微服务冷启动可能是秒级，而大模型服务可能需要几十秒到数分钟。冷启动包括：

\begin{itemize}
  \item 镜像拉取；
  \item 模型权重从对象存储或分布式文件系统加载；
  \item 权重反序列化和分片；
  \item GPU 显存分配；
  \item CUDA context 初始化；
  \item NCCL communicator 初始化；
  \item kernel autotuning；
  \item CUDA graph capture；
  \item KV Cache block 初始化；
  \item readiness probe 通过。
\end{itemize}

降低冷启动的策略包括：

\begin{itemize}
  \item 使用本地 NVMe 缓存模型权重；
  \item 预拉取镜像；
  \item 使用 warm pool；
  \item 模型分层加载；
  \item 权重格式优化；
  \item 避免频繁 scale-to-zero；
  \item 将热点模型常驻；
  \item 使用灰度流量切换，避免未预热实例接收真实请求。
\end{itemize}

\subsection{多模型共存}
\label{subsec:multi-model-coexistence}

一个生产集群通常同时服务多个模型：大模型、小模型、embedding、rerank、多模态、代码模型、专用领域模型、实验版本和灰度版本。多模型共存带来几个调度问题：

\begin{itemize}
  \item 热点模型需要更多副本，冷门模型是否常驻？
  \item 多个小模型是否合并到同一 serving runtime？
  \item LoRA adapter 是否动态加载？
  \item 模型之间如何共享 KV Cache 或 prefix cache？
  \item 不同模型的 SLA 是否不同？
  \item 模型升级时如何做流量切换和回滚？
\end{itemize}

\begin{figure}[H]
  \centering
  \begin{tikzpicture}[
    font=\small,
    box/.style={draw, rounded corners, minimum width=2.1cm, minimum height=0.7cm, align=center},
    wide/.style={draw, rounded corners, minimum width=8.5cm, minimum height=0.75cm, align=center},
    arrow/.style={-{Stealth[length=2mm]}, thick}
  ]
    \node[wide] (client) at (0,5.0) {Clients / Applications};
    \node[wide] (gw) at (0,4.0) {API Gateway / Auth / Rate Limit};
    \node[wide] (router) at (0,3.0) {Model Router / Admission Control};

    \node[box] (s0) at (-3.2,1.8) {Prefill Pool\\GPU Pods};
    \node[box] (s1) at (0,1.8) {Decode Pool\\GPU Pods};
    \node[box] (s2) at (3.2,1.8) {Embedding /\\Rerank};

    \node[wide] (obs) at (0,0.6) {Metrics / Logs / Traces / Cost Accounting};
    \node[wide] (store) at (0,-0.55) {Model Registry / Object Storage / NVMe Cache};

    \draw[arrow] (client) -- (gw);
    \draw[arrow] (gw) -- (router);
    \draw[arrow] (router) -- (s0);
    \draw[arrow] (router) -- (s1);
    \draw[arrow] (router) -- (s2);
    \draw[arrow] (s0) -- (obs);
    \draw[arrow] (s1) -- (obs);
    \draw[arrow] (s2) -- (obs);
    \draw[arrow] (store) -- (s0);
    \draw[arrow] (store) -- (s1);
    \draw[arrow] (store) -- (s2);
  \end{tikzpicture}
  \caption{LLM Serving on Kubernetes 架构图：模型路由、GPU Pod、观测系统和模型存储共同构成在线服务平台。}
  \label{fig:llm-serving-on-k8s}
\end{figure}

\section{AI 集群可观测性}
\label{sec:ai-cluster-observability}

没有可观测性，就没有真正的调度优化。AI 集群的可观测性必须跨越 Kubernetes、GPU、网络、存储、作业和业务指标。只看 GPU utilization 会误导决策；只看 Pod 状态也无法解释训练慢在哪里。

\subsection{GPU Utilization}
\label{subsec:gpu-utilization}

GPU utilization 是最常见指标，但也是最容易被误解的指标。它通常表示一段采样时间内 GPU 是否有活跃 work，而不是 Tensor Core 是否高效运行。一个 NCCL kernel、一个 memory-bound kernel、一个低效小 kernel 都可能让 GPU utilization 看起来不低，但业务吞吐并不好。

因此，GPU utilization 应与以下指标一起看：

\begin{itemize}
  \item SM utilization；
  \item Tensor Core utilization；
  \item HBM bandwidth；
  \item GPU memory usage；
  \item kernel timeline；
  \item NCCL communication time；
  \item tokens per second 或 samples per second；
  \item cost per token 或 cost per step。
\end{itemize}

\subsection{SM Occupancy}
\label{subsec:sm-occupancy-monitoring}

SM occupancy 或 SM active 反映计算单元活跃程度。训练中的大 GEMM 通常希望 SM 和 Tensor Core 利用率较高；推理 decode 阶段则可能因为 batch 小、KV Cache 读取和小 GEMM 导致 SM 利用率不高。

调度系统可以利用 SM 指标判断混部机会。例如一个服务显存占用高但 SM 使用低，理论上可以与另一个计算型任务混部；但如果它对延迟敏感，混部可能导致 SLA 抖动。因此，资源利用率指标必须与 workload 类型和 SLA 结合。

\subsection{HBM Usage}
\label{subsec:hbm-usage-monitoring}

HBM usage 是训练和推理都非常关键的指标。训练中，HBM 被参数、梯度、优化器状态、activation 和临时 buffer 占用；推理中，HBM 被模型权重、KV Cache、workspace 和 runtime buffer 占用。

推理服务尤其需要监控 KV Cache 使用率。如果 KV Cache block 使用率接近上限，新的请求可能被拒绝、排队或触发 eviction。调度系统可以根据 KV Cache 压力而不是只根据 QPS 做扩容。

\subsection{Network Throughput}
\label{subsec:network-throughput-monitoring}

网络指标包括节点内和节点间两类：

\begin{itemize}
  \item 节点内：NVLink/NVSwitch 带宽、PCIe 带宽、GPU P2P；
  \item 节点间：NIC TX/RX、RDMA counters、PFC、ECN、丢包、重传、交换机端口利用率。
\end{itemize}

对于训练作业，需要把 NCCL collective 时间与网络指标关联起来。如果某个 job step time 变长，同时某些 leaf-spine 链路利用率接近上限，说明可能存在网络热点。如果 NCCL timeout 与 PFC pause storm 同时出现，说明 RoCE 拥塞控制需要排查。

\subsection{Job Queue Time}
\label{subsec:job-queue-time}

队列等待时间是 AI 平台用户最关心的指标之一。一个训练任务真正的完成时间不是运行时间，而是：

\[
  T_{\text{completion}} = T_{\text{queue}} + T_{\text{startup}} + T_{\text{run}} + T_{\text{recovery}}.
\]

其中 $T_{\text{queue}}$ 可能非常长。大作业经常因为资源碎片等待，小作业则可能因为低优先级等待。平台应监控：

\begin{itemize}
  \item 不同队列的平均等待时间；
  \item P95/P99 等待时间；
  \item 不同 GPU flavor 的等待时间；
  \item 大作业与小作业等待差异；
  \item 被抢占次数；
  \item backfill 成功率；
  \item 资源碎片率。
\end{itemize}

队列时间可以反向指导容量规划。如果 H100 队列长期拥塞，而 A100 队列空闲，说明资源结构与需求不匹配；如果 8 GPU 单节点作业长期等待，但单卡任务很快启动，说明 bin packing 或 reservation 策略需要调整。

\subsection{Cost Per Token / Cost Per Step}
\label{subsec:cost-per-token-step}

最终，AI Infra 的目标不是让某个指标好看，而是降低单位业务产出的成本。训练中常用 cost per step、cost per token、tokens per GPU-hour；推理中常用 cost per 1K tokens、cost per request、GPU-hour per SLA bucket。

这些指标把硬件利用率、调度效率和业务结果连接起来。例如：

\begin{itemize}
  \item 一个推理服务 GPU utilization 很高，但 P99 延迟也很高，可能需要扩容；
  \item 一个训练作业 GPU utilization 不低，但 tokens per GPU-hour 低，可能是通信瓶颈；
  \item 一个队列 GPU 使用率高，但大量作业被抢占，实际有效训练 token 低；
  \item 一个量化模型成本低，但质量下降导致重试率上升，端到端成本未必低。
\end{itemize}

\begin{figure}[H]
  \centering
  \begin{tikzpicture}[
    font=\small,
    layer/.style={draw, rounded corners, minimum width=9.0cm, minimum height=0.7cm, align=center},
    arrow/.style={-{Stealth[length=2mm]}, thick}
  ]
    \node[layer] (biz) at (0,5.0) {业务层：QPS、Latency、SLA、Cost per Token、用户体验};
    \node[layer] (job) at (0,4.0) {作业层：Queue Time、Run Time、Checkpoint、Failure、Preemption};
    \node[layer] (k8s) at (0,3.0) {Kubernetes 层：Pod、Node、Scheduler、Quota、Events};
    \node[layer] (gpu) at (0,2.0) {GPU 层：Utilization、SM、HBM、Tensor Core、Xid、ECC};
    \node[layer] (net) at (0,1.0) {网络层：NVLink、RDMA、NIC、PFC、ECN、NCCL Bandwidth};
    \node[layer] (storage) at (0,0.0) {存储层：Dataset、Checkpoint、Model Loading、NVMe Cache};

    \draw[arrow] (storage) -- (net);
    \draw[arrow] (net) -- (gpu);
    \draw[arrow] (gpu) -- (k8s);
    \draw[arrow] (k8s) -- (job);
    \draw[arrow] (job) -- (biz);
  \end{tikzpicture}
  \caption{AI Infra 监控指标分层图：从硬件资源到业务成本需要跨层关联。}
  \label{fig:ai-infra-monitoring-layers}
\end{figure}

\section{从 Kubernetes 到 AI 平台}
\label{sec:from-kubernetes-to-ai-platform}

Kubernetes 是强大的基础设施底座，但一个完整 AI 平台通常还需要更上层的能力：

\begin{itemize}
  \item 训练任务提交系统；
  \item 模型注册与版本管理；
  \item 数据集管理；
  \item 实验追踪；
  \item 分布式 checkpoint 管理；
  \item 镜像构建与依赖管理；
  \item 权限、审计和成本分摊；
  \item 自动调参和超参搜索；
  \item 模型评测流水线；
  \item 推理流量治理；
  \item 多集群和多区域调度；
  \item 故障自动诊断。
\end{itemize}

因此，不要把“在 Kubernetes 上跑 GPU Pod”误认为“已经有了 AI 平台”。真正的平台需要把用户体验、资源效率、可靠性和成本闭环起来。对用户来说，他们希望提交一个训练配置或模型部署配置，而不是手写复杂 YAML；对平台来说，需要把这些高级意图翻译成队列、配额、Pod、Service、PVC、Secret、ConfigMap、Job、监控和告警。

一个成熟 AI 平台通常会形成如下分层：

\begin{enumerate}
  \item \textbf{硬件层}：GPU、NIC、NVMe、交换机、机架和供电；
  \item \textbf{系统层}：驱动、CUDA、container runtime、RDMA、存储客户端；
  \item \textbf{Kubernetes 层}：Pod、Node、Service、Scheduler、Device Plugin、Operator；
  \item \textbf{Batch / Serving 层}：Volcano、Kueue、Training Operator、Serving Runtime；
  \item \textbf{平台层}：任务提交、模型管理、数据管理、权限、计费、可观测性；
  \item \textbf{业务层}：训练产出、推理 SLA、模型质量、成本收益。
\end{enumerate}

AI Infra 工程师的价值就在于能跨越这些层次定位问题。例如用户反馈“训练慢”，原因可能是模型代码、GEMM shape、NCCL all-reduce、RoCE 拥塞、GPU-NIC 亲和、Pod 被错误调度、存储读取慢或 checkpoint 阻塞。只懂其中一层，很难完成端到端优化。

\section{本章小结}
\label{sec:k8s-ai-scheduling-summary}

本章介绍了 Kubernetes 在 AI 集群中的调度角色。我们从 GPU 资源昂贵、多租户、训练推理混部、利用率与公平性的矛盾出发，解释了为什么 AI 集群需要比普通 Kubernetes 更强的调度能力。

本章核心要点如下：

\begin{itemize}
  \item AI 集群调度不是简单分配 GPU 数量，而是要同时考虑 GPU 型号、显存、拓扑、网络、队列、优先级、SLA 和成本。
  \item Kubernetes 的 Pod、Node、Scheduler、Device Plugin 和 Operator 构成了 GPU 集群管理的基础。
  \item NVIDIA Device Plugin 把 GPU 暴露为 Kubernetes 扩展资源，GPU Operator 可以进一步管理驱动、监控、MIG 等组件。
  \item MIG 适合小模型推理和多租户隔离，GPU sharing 适合低优先级或可容忍抖动的共享场景，但二者都需要谨慎处理性能隔离。
  \item 拓扑感知调度对 TP、EP、多 rail RDMA 和多卡推理至关重要，不能只按 GPU count 调度。
  \item 分布式训练需要 gang scheduling，避免 worker 部分启动造成 GPU 浪费。
  \item Volcano 提供面向 batch/HPC/AI workload 的队列、gang、优先级、抢占和 backfill 等能力。
  \item Kueue 提供 Kubernetes 原生作业队列、配额和准入控制，适合多租户 batch/AI 集群。
  \item 推理服务编排关注 Deployment、autoscaling、bin packing、cold start、多模型共存和 SLA。
  \item AI 集群可观测性必须跨越业务、作业、Kubernetes、GPU、网络和存储层，最终落到 cost per token 与 cost per step。
\end{itemize}

至此，本书第五篇完成了从 GPU 硬件到集群调度的完整路径：单卡计算、单机互联、多机网络、Kubernetes 调度。结合前面几篇内容，我们已经从数学、CUDA、Transformer、分布式训练、推理优化一路走到了 AI Infrastructure 的完整系统视角。后续附录将补充符号约定、CUDA 与分布式环境搭建、AI 集群上线检查清单以及参考资料索引，帮助读者把本书内容进一步落到工程实践中。
